\chapter[35]{Stability of Ricci Flow}
\label{ch:chow-iv-stability-ricci-flow}
\dcref{ch35}
\source{chow-et-al-ricci-flow-part-iv:page-0300}

\section{Linear stability of Ricci flow}

For a parabolic system
\begin{equation*}
 \partial_tu=A(x,t,u,Du,D^2u),\qquad u(x,0)=u_0(x),
 \tag{35.1}
\end{equation*}
linearization at a fixed point determines the first-order stable and unstable
directions.

\begin{definition}[Linear stability (Definition 35.1)]
\label{def:chow-iv-linear-stability}
\source{chow-et-al-ricci-flow-part-iv:page-0301}
\dcref{ch35:35.1}
\group{Linear Stability}

Let $\Sigma$ be the spectrum of the elliptic linearization
$D_{\bar u}A$ and let $\sigma$ be its projection to $\mathbb R$.  The fixed
point is strictly linearly stable if $\sigma\subset(-\infty,0)$, linearly
stable if $\sigma\subset(-\infty,0]$, and linearly unstable if
$\sigma\cap(0,\infty)\ne\varnothing$.
\end{definition}

Let $\mathcal M^n$ be closed and connected.  Write $S^2$ and $S^2_+$ for the
spaces of smooth symmetric and positive-definite $(2,0)$-tensors.  The
divergence and Killing operators are
\[
 \delta h=-g^{ij}\nabla_i h_{jk}\,dx^k,
 \qquad \delta^*\omega=\tfrac12\mathcal L_{\omega^\#}g.
\]
Ebin's slice decomposition gives
\begin{equation*}
 T_gS^2_+=\mathcal H_g\oplus\mathcal V_g,
 \qquad \mathcal H_g=\ker\delta_g,\quad
 \mathcal V_g=\operatorname{im}\delta_g^*.
 \tag{35.2}
\end{equation*}

For a background metric $\hat g$, set
\begin{equation*}
 A_{\hat g}(g)=-2\Rc(g)-P_{\hat g}(g),
 \qquad P_{\hat g}(g)_{ij}=-(\nabla_iW_j+\nabla_jW_i),
 \tag{35.3}
\end{equation*}
where
\[
 W_j=(\hat g^{-1})_{jk}g^{k\ell}g^{pq}
 \left(\nabla_pg_{q\ell}-\tfrac12\nabla_\ell g_{pq}\right).
\]
The Ricci--DeTurck flow is
\begin{equation*}
 \partial_tg=A_{\hat g}(g),\qquad g(0)=g_0.
 \tag{35.5}
\end{equation*}

\begin{proposition}[Lichnerowicz linearization (Proposition 35.3)]
\label{prop:chow-iv-lichnerowicz-linearization}
\source{chow-et-al-ricci-flow-part-iv:page-0304}
\dcref{ch35:35.3}
\group{Linear Stability}

If $\hat g=\bar g$ is a Ricci-flow fixed point, then
\[
 (D_hA_{\bar g})(\bar g)=\Delta_Lh,
 \qquad
 (\Delta_Lh)_{ij}=\Delta h_{ij}+2R_{kij\ell}h_{k\ell}-R_{ik}h_{kj}-R_{jk}h_{ki}.
 \tag{35.7}
\]
The operator $\Delta_L$ is self-adjoint and strictly elliptic on a compact
manifold.
\end{proposition}

Shrinking and expanding solitons become stationary under the dilated flows
\[
 \partial_\tau\gamma=-2\Rc(\gamma)-\mathcal L_X\gamma+2\lambda\gamma,
 \qquad
 \partial_\tau\hat g=-2\Rc(\hat g)-\mathcal L_X\hat g-2\lambda\hat g.
\]

\begin{proposition}[Flat metrics (Proposition 35.5)]
\label{prop:chow-iv-flat-linear-stability}
\source{chow-et-al-ricci-flow-part-iv:page-0305}
\dcref{ch35:35.5}
\group{Linear Stability}

If $(\mathcal M^n,g)$ is flat, then $\Delta_L=\Delta$ is negative
semidefinite on $S^2$, with kernel exactly the parallel $(2,0)$-tensors.
Thus the kernel has dimension at most $n(n+1)/2$ and every flat metric is
linearly stable for Ricci--DeTurck flow.
\end{proposition}

\begin{definition}[$K3$ surface (Definition 35.6)]
\label{def:chow-iv-k3-surface}
\source{chow-et-al-ricci-flow-part-iv:page-0305}
\dcref{ch35:35.6}
\group{Linear Stability}

A $K3$ surface is a closed connected smooth complex surface with vanishing
first Chern class and no global holomorphic $1$-form.
\end{definition}

\begin{proposition}[Linear stability of $K3$ surfaces (Proposition 35.7)]
\label{prop:chow-iv-k3-linear-stability}
\source{chow-et-al-ricci-flow-part-iv:page-0305}
\uses{def:chow-iv-k3-surface}
\dcref{ch35:35.7}
\group{Linear Stability}


Every K\"ahler--Einstein $K3$ surface is linearly stable for Ricci--DeTurck
flow.  The kernel of $\Delta_L$ is $\varepsilon(g)\oplus G$, where
$\varepsilon(g)$ is the product of the $3$-dimensional space of parallel
self-dual $2$-forms and the $19$-dimensional space of harmonic anti-self-dual
$2$-forms.
\end{proposition}

\begin{theorem}[Parallel-spinor manifolds (Theorem 35.8)]
\label{thm:chow-iv-parallel-spinor-stability}
\source{chow-et-al-ricci-flow-part-iv:page-0305}
\dcref{ch35:35.8}
\group{Linear Stability}

If a compact Riemannian manifold is covered by a spin manifold admitting a
nonzero parallel spinor, then $\Delta_L$ has nonpositive spectrum.
\end{theorem}

\begin{remark}[Parallel-spinor consequence (Remark 35.9)]
\label{rem:chow-iv-parallel-spinor-rigidity}
\source{chow-et-al-ricci-flow-part-iv:page-0305}
\uses{thm:chow-iv-parallel-spinor-stability}
\dcref{ch35:35.9}
\group{Linear Stability}


The hypothesis forces Ricci flatness and generalizes the $K3$ result.
\end{remark}

\begin{proposition}[Ricci-flat holonomy (Proposition 35.10)]
\label{prop:chow-iv-ricci-flat-holonomy}
\source{chow-et-al-ricci-flow-part-iv:page-0306}
\dcref{ch35:35.10}
\group{Linear Stability}

A simply-connected irreducible compact Ricci-flat manifold either has holonomy
$\mathrm{SO}(n)$ or admits a nonzero parallel spinor.
\end{proposition}

\begin{remark}[Known compact Ricci-flat examples (Remark 35.11)]
\label{rem:chow-iv-known-ricci-flat-stability}
\source{chow-et-al-ricci-flow-part-iv:page-0306}
\dcref{ch35:35.11}
\group{Linear Stability}

All known compact Ricci-flat examples have special holonomy and hence are
linearly stable by the parallel-spinor theorem.
\end{remark}

\begin{theorem}[Parallel $\mathrm{spin}^c$ spinors (Theorem 35.12)]
\label{thm:chow-iv-spinc-linear-stability}
\source{chow-et-al-ricci-flow-part-iv:page-0306}
\dcref{ch35:35.12}
\group{Linear Stability}

If a compact Einstein manifold with nonpositive scalar curvature admits a
nonzero parallel $\mathrm{spin}^c$ spinor, then
\[
 Ah_{ij}=\Delta_Lh_{ij}+R_i{}^kh_{kj}+R_j{}^kh_{ik}
\]
has nonpositive spectrum.  In particular this holds for compact
K\"ahler--Einstein manifolds with nonpositive scalar curvature.
\end{theorem}

An algebraic soliton on a homogeneous space is a left-invariant metric satisfying
\[
 \Rc=\lambda I+D,
\]
with $D$ a derivation.  Such expanding solitons include all nilsolitons and
many solvsolitons; nilsolitons of dimensions at most six and several
solvsoliton families are strictly linearly stable.

\begin{proposition}[Second variation of $\lambda$ (Proposition 35.13)]
\label{prop:chow-iv-perelman-lambda-second-variation}
\source{chow-et-al-ricci-flow-part-iv:page-0307,chow-et-al-ricci-flow-part-iv:page-0308}
\dcref{ch35:35.13}
\group{Linear Stability}

For a compact Ricci-flat manifold, Perelman's average energy
\[
 \lambda(g)=\inf_{\int e^{-f}d\mu=1}\int(|\nabla f|^2+R)e^{-f}d\mu
 \tag{35.8}
\]
has second variation
\[
 D^2\lambda(g)(h,h)=\int\langle Lh,h\rangle d\mu,
 \qquad
 Lh=\begin{cases}\frac12\Delta_Lh,&h\in\mathcal H,\\0,&h\in\mathcal V.
 \end{cases}
\]
\end{proposition}

\begin{definition}[Dynamic stability (Definition 35.14)]
\label{def:chow-iv-dynamic-stability}
\source{chow-et-al-ricci-flow-part-iv:page-0308}
\uses{def:chow-iv-linear-stability}
\dcref{ch35:35.14}
\group{Dynamic Stability}


For $\partial_tg=F(g)$ with fixed point $\bar g$, asymptotic stability means
that all nearby solutions exist for positive time and converge to $\bar g$.
Stability means that every neighborhood $V$ of $\bar g$ contains a neighborhood
$U$ such that solutions starting in $U$ remain in $V$ for all positive time.
\end{definition}

\section{Analytic semigroups and maximal regularity}

For Banach spaces $\mathcal E_1\hookrightarrow\mathcal E_0$ and a densely
defined operator $A$, the linear problem is
\begin{equation*}
 u'=Au+f,\qquad u(0)=0,
 \tag{35.10}
\end{equation*}
with formal solution
\begin{equation*}
 u(t)=Sf(t)=\int_0^t e^{(t-s)A}f(s)\,ds.
 \tag{35.11}
\end{equation*}
An operator is sectorial when its resolvent exists in a right half-plane and
satisfies the standard bound $\|R(\lambda,A)f\|_{\mathcal E_1}\le
C|\lambda-\omega|^{-1}\|f\|_{\mathcal E_0}$; then $e^{tA}$ is an analytic
semigroup.

\begin{proposition}[Linear maximal-regularity estimate (Proposition 35.18)]
\label{prop:chow-iv-linear-maximal-regularity}
\source{chow-et-al-ricci-flow-part-iv:page-0311}
\dcref{ch35:35.18}
\group{Maximal Regularity}

If $f\in C([0,T];\mathcal E_\theta)$ and the spectral bound
$\omega=\sup\{\operatorname{Re}\lambda:\lambda\in\sigma(A)\}<0$, then for
$\eta+\omega<0$,
\begin{equation*}
 \sup_t\|e^{\eta t}u'(t)\|_{\mathcal E_\theta}
 +\sup_t\|e^{\eta t}u(t)\|_{\mathcal E_{1+\theta}}
 \le C\sup_t\|e^{\eta t}f(t)\|_{\mathcal E_\theta}.
 \tag{35.12}
\end{equation*}
\end{proposition}

\begin{proposition}[Nonlinear asymptotic stability estimate (Proposition 35.19)]
\label{prop:chow-iv-nonlinear-maximal-regularity}
\source{chow-et-al-ricci-flow-part-iv:page-0311,chow-et-al-ricci-flow-part-iv:page-0312}
\uses{prop:chow-iv-linear-maximal-regularity}
\dcref{ch35:35.19}
\group{Maximal Regularity}


For $u'=F(u)$ with $F(0)=0$, sectorial linear part, $C^1$ nonlinearity,
and locally Lipschitz derivative, there is $r>0$ such that
$\|u_0\|_{\mathcal E_{1+\theta}}\le r$ gives a unique exponentially decaying
solution
\[
 u\in C^1_\eta([0,\infty);\mathcal E_\theta)\cap
 C_\eta([0,\infty);\mathcal E_{1+\theta}).
\]
\end{proposition}

\subsection{Little H\"older spaces}

\begin{definition}[Little H\"older spaces (Definitions 35.20--35.21)]
\label{def:chow-iv-little-holder-spaces}
\source{chow-et-al-ricci-flow-part-iv:page-0313}
\dcref{ch35:35.20--35.21}
\group{Maximal Regularity}

The little H\"older space $h^{r+\rho}$ is the closure of smooth functions (or
smooth symmetric $2$-tensors on $\mathcal M$) in the $C^{r,\rho}$ norm.  For
$s\le r$ and $0<\sigma<\rho<1$ the inclusion
$h^{r+\rho}\hookrightarrow h^{s+\sigma}$ is continuous and dense.
\end{definition}

\begin{definition}[Continuous interpolation (Definition 35.22)]
\label{def:chow-iv-continuous-interpolation}
\source{chow-et-al-ricci-flow-part-iv:page-0314}
\dcref{ch35:35.22}
\group{Maximal Regularity}

For Banach spaces $B_1\subset B_0$ and $0<\theta<1$, $(B_0,B_1)_\theta$
consists of those $x=y_n+z_n$ with
$\|y_n\|_{B_0}=o(2^{-n\theta})$ and
$\|z_n\|_{B_1}=o(2^{n(1-\theta)})$; its norm is the corresponding infimum of
$\sup_n\max(2^{n\theta}\|y_n\|,2^{-n(1-\theta)}\|z_n\|)$.
\end{definition}

\begin{lemma}[Interpolation of little H\"older spaces (Lemma 35.24)]
\label{lem:chow-iv-holder-interpolation}
\source{chow-et-al-ricci-flow-part-iv:page-0314}
\dcref{ch35:35.24}
\group{Maximal Regularity}

If the interpolated exponent is not an integer, then
\begin{equation*}
 (h^{s+\sigma},h^{r+\rho})_\theta\cong
 h^{(\theta r+(1-\theta)s)+\theta\rho+(1-\theta)\sigma},
 \tag{35.15}
\end{equation*}
and
\begin{equation*}
 \|\eta\|_{(h^{s+\sigma},h^{r+\rho})_\theta}
 \le C\|\eta\|_{h^{s+\sigma}}^{1-\theta}
 \|\eta\|_{h^{r+\rho}}^\theta.
 \tag{35.16}
\end{equation*}
\end{lemma}

Set $\mathcal E_0=h^{0+\sigma}$, $\mathcal X_0=h^{0+\rho}$,
$\mathcal E_1=h^{2+\sigma}$, and $\mathcal X_1=h^{2+\rho}$, with
$\theta=(\rho-\sigma)/2$ so that
\begin{equation*}
 \mathcal X_0\cong(\mathcal E_0,\mathcal E_1)_\theta,
 \qquad \mathcal X_1\cong(\mathcal E_0,\mathcal E_1)_{1+\theta}.
 \tag{35.18}
\end{equation*}

\begin{lemma}[Ricci--DeTurck analytic dependence (Lemma 35.25)]
\label{lem:chow-iv-deturck-analytic-dependence}
\source{chow-et-al-ricci-flow-part-iv:page-0315}
\dcref{ch35:35.25}
\group{Maximal Regularity}

For the positive-metric neighborhoods $\mathcal G_\beta$ and $\mathcal G_\alpha$
defined from these spaces, $g\mapsto A(g)$ and
$g\mapsto\widetilde{\mathcal A}(g)$ are analytic maps into
$\mathcal L(\mathcal X_1,\mathcal X_0)$ and
$\mathcal L(\mathcal E_1,\mathcal E_0)$, respectively.
\end{lemma}

\begin{lemma}[Analytic semigroup for Ricci--DeTurck (Lemma 35.26)]
\label{lem:chow-iv-deturck-analytic-semigroup}
\source{chow-et-al-ricci-flow-part-iv:page-0315,chow-et-al-ricci-flow-part-iv:page-0316}
\uses{lem:chow-iv-deturck-analytic-dependence}
\dcref{ch35:35.26}
\group{Maximal Regularity}


For every $g\in\mathcal G_\beta$, the strongly elliptic extension
$\widetilde{\mathcal A}(g):\mathcal E_1\subset\mathcal E_0\to\mathcal E_0$
generates a strongly continuous analytic semigroup.
\end{lemma}

\begin{theorem}[Center Manifold Theorem (Theorem 35.27)]
\label{thm:chow-iv-center-manifold}
\source{chow-et-al-ricci-flow-part-iv:page-0316,chow-et-al-ricci-flow-part-iv:page-0317}
\uses{lem:chow-iv-holder-interpolation, lem:chow-iv-deturck-analytic-semigroup}
\dcref{ch35:35.27}
\group{Maximal Regularity}


For an autonomous quasilinear parabolic equation
\begin{equation*}
 \partial_tg=A(g)g,
 \tag{35.19}
\end{equation*}
whose linearization at a fixed point has stable spectrum and finitely many
center/unstable eigenvalues, the interpolation spaces split into stable and
center parts.  For every $r$ there is a local $C^r$ graph
$\mathcal M^c_{\mathrm{loc}}$ tangent to the center eigenspace, locally
invariant while solutions remain nearby, and stable components approach the
graph at rate $t^{\alpha-1}e^{-\omega t}$.  The theorem gives no uniqueness
of the graph in general and no dynamics inside its center part.
\end{theorem}

\section{Dynamic stability obtained by linearization}

\begin{lemma}[Convergence of DeTurck diffeomorphisms (Lemma 35.28)]
\label{lem:chow-iv-deturck-diffeomorphism-convergence}
\source{chow-et-al-ricci-flow-part-iv:page-0318}
\dcref{ch35:35.28}
\group{Dynamic Stability}

If a time-dependent vector field satisfies
$\sup_x|W(x,t)|_{g(t)}\le Ce^{-ct}$, then its generated diffeomorphisms
converge exponentially to a limiting diffeomorphism.
\end{lemma}

\begin{corollary}[From Ricci--DeTurck to Ricci flow (Corollary 35.29)]
\label{cor:chow-iv-deturck-to-ricci-flow}
\source{chow-et-al-ricci-flow-part-iv:page-0318}
\uses{lem:chow-iv-deturck-diffeomorphism-convergence}
\dcref{ch35:35.29}
\group{Dynamic Stability}


If Ricci--DeTurck flow near a flat metric converges exponentially to a flat
metric, then Ricci flow with the same initial data also converges exponentially
to a flat metric.
\end{corollary}

\begin{theorem}[Asymptotic stability at $\Rm=0$ (Theorem 35.30)]
\label{thm:chow-iv-flat-asymptotic-stability}
\source{chow-et-al-ricci-flow-part-iv:page-0319,chow-et-al-ricci-flow-part-iv:page-0320}
\uses{thm:chow-iv-center-manifold, cor:chow-iv-deturck-to-ricci-flow}
\dcref{ch35:35.30}
\group{Dynamic Stability}


For a flat compact $(\mathcal M^n,\bar g)$, the tangent space splits into a
stable part and the $n(n+1)/2$-dimensional space of parallel tensors.  There
is a unique smooth exponentially attractive center manifold consisting exactly
of nearby flat metrics, and every Ricci-flow solution starting sufficiently
near $\bar g$ converges exponentially to a flat metric on it.
\end{theorem}

\subsection{Dynamic stability in negative curvature}

For an unnormalized flow existing on $(T_0,T_1)$, set
$\sigma(t)=2r(t-T_0)$ and $\tau=(2r)^{-1}\log(t-T_0)$.  The rescaled metric
satisfies the dilated flow
\begin{equation*}
 \partial_\tau\widetilde g=-2\Rc-2r\widetilde g.
 \tag{35.21}
\end{equation*}
At an Einstein metric with $\Rc=-rg$, the DeTurck linearization is
$Lh=\Delta_Lh-2rh=\Delta h+2R_{ipqj}h^{pq}$.  With
$Q(h)=R_{ijkl}h^{i\ell}h^{jk}$ and
$T_{ijk}=\nabla_kh_{ij}-\nabla_ih_{jk}$,
\begin{align*}
 (Lh,h)&=-\|\nabla h\|^2+2\int Q(h)\,d\mu, \tag{35.22}\\
 (Lh,h)&=-\tfrac12\|T\|^2-\|\delta h\|^2-r\|h\|^2+\int Q(h)\,d\mu. \tag{35.23}
\end{align*}
If $W$ is the Weyl tensor, then
\begin{equation*}
 (Lh,h)\le -\frac{n-2}{n-1}r\|h\|^2+\int W(h,h)\,d\mu.
 \tag{35.24}
\end{equation*}

\begin{theorem}[Asymptotic stability of negative curvature (Theorem 35.31)]
\label{thm:chow-iv-negative-curvature-stability}
\source{chow-et-al-ricci-flow-part-iv:page-0321,chow-et-al-ricci-flow-part-iv:page-0322}
\uses{prop:chow-iv-nonlinear-maximal-regularity}
\dcref{ch35:35.31}
\group{Dynamic Stability}


If $n>2$ and $(\mathcal M^n,\bar g)$ is compact with constant sectional
curvature $K<0$, then normalized Ricci flow from a sufficiently small
$C^{2+\alpha}$ little H\"older neighborhood exists for all time and converges
exponentially to a constant-curvature metric $\lambda\bar g$.
\end{theorem}

\begin{remark}[Regularity boost (Remark 35.32)]
\label{rem:chow-iv-negative-curvature-regularity}
\source{chow-et-al-ricci-flow-part-iv:page-0322}
\dcref{ch35:35.32}
\group{Dynamic Stability}

Initial data in a $1+\beta$ little H\"older neighborhood converge in the
$2+\alpha$ topology for suitable $\beta>\alpha$.
\end{remark}

\begin{remark}[Curvature pinching is insufficient (Remark 35.33)]
\label{rem:chow-iv-negative-curvature-obstruction}
\source{chow-et-al-ricci-flow-part-iv:page-0322}
\dcref{ch35:35.33}
\group{Dynamic Stability}

For dimensions $n>10$, there are hyperbolic manifolds with metrics whose
sectional curvatures are arbitrarily pinched near $-1$ but whose Ricci flows
do not converge smoothly to the hyperbolic metric; the theorem requires
topological closeness in a H\"older norm.
\end{remark}

\section{Further dynamic stability results}

For normalized Ricci--DeTurck flow at an Einstein metric, the linearization is
\begin{equation*}
 Ah=\Delta_Lh+2K\left(h-\frac1n\bar H g\right),
 \qquad \bar H=\frac{\int\operatorname{tr}_gh\,d\mu}{\int d\mu}.
 \tag{35.26}
\end{equation*}
At positive constant curvature, the only unstable eigendirections are the
conformal-diffeomorphism directions on the round sphere (and are therefore
ungeometric).  Cross-curvature flow, renormalization-group flow, and stability
results for collapsed limits provide related applications of maximal regularity.

\begin{remark}[Dynamic stability of positive curvature (Remark 35.34)]
\label{rem:chow-iv-positive-curvature-stability}
\source{chow-et-al-ricci-flow-part-iv:page-0325}
\dcref{ch35:35.34}
\group{Alternative Stability Methods}

Simple connectivity is imposed below to ensure uniqueness of the limiting
space form up to isometry.
\end{remark}

\begin{theorem}[Hamilton in dimension three (Theorem 35.35)]
\label{thm:chow-iv-hamilton-positive-ricci}
\source{chow-et-al-ricci-flow-part-iv:page-0325}
\dcref{ch35:35.35}
\group{Alternative Stability Methods}

If a closed simply-connected $3$-manifold has positive Ricci curvature, its
volume-normalized Ricci flow converges to a unique metric of constant positive
curvature.
\end{theorem}

\begin{theorem}[Hamilton in dimension four (Theorem 35.36)]
\label{thm:chow-iv-hamilton-positive-curvature-operator}
\source{chow-et-al-ricci-flow-part-iv:page-0325}
\dcref{ch35:35.36}
\group{Alternative Stability Methods}

If a closed simply-connected $4$-manifold has positive curvature operator, its
volume-normalized Ricci flow converges to a unique constant-positive-curvature
metric; the manifold is the standard sphere (or projective space if
nonorientable).
\end{theorem}

\begin{theorem}[Huisken pinching (Theorem 35.37)]
\label{thm:chow-iv-huisken-pinching}
\source{chow-et-al-ricci-flow-part-iv:page-0325,chow-et-al-ricci-flow-part-iv:page-0326}
\dcref{ch35:35.37}
\group{Alternative Stability Methods}

If a closed simply-connected $n$-manifold, $n\ge4$, has positive scalar
curvature and irreducible curvature components satisfying
\[
 |V|^2+|W|^2<\delta_n|U|^2=\frac{2\delta_n}{n(n-1)}R^2,
\]
then volume-normalized Ricci flow converges to a unique constant-positive-
curvature metric.
\end{theorem}

\begin{remark}[Pointwise pinching (Remark 35.38)]
\label{rem:chow-iv-pointwise-pinching}
\source{chow-et-al-ricci-flow-part-iv:page-0326}
\dcref{ch35:35.38}
\group{Alternative Stability Methods}

The pinching hypothesis is pointwise rather than a global sectional-curvature
pinching condition.
\end{remark}

\begin{theorem}[Margerin's optimal four-dimensional pinching (Theorem 35.39)]
\label{thm:chow-iv-margerin-pinching}
\source{chow-et-al-ricci-flow-part-iv:page-0326}
\dcref{ch35:35.39}
\group{Alternative Stability Methods}

For a closed simply-connected $4$-manifold with positive scalar curvature,
$|V|^2+|W|^2<|U|^2=R^2/6$ implies convergence to a unique constant-positive-
curvature metric.  Equality yields the Fubini--Study metric on $\mathbb{CP}^2$
or a quotient of $\mathbb R\times S^3$.
\end{theorem}

\begin{theorem}[Two-positive curvature operator (Theorem 35.40)]
\label{thm:chow-iv-two-positive-curvature}
\source{chow-et-al-ricci-flow-part-iv:page-0326}
\dcref{ch35:35.40}
\group{Alternative Stability Methods}

If a closed simply-connected $n$-manifold, $n\ge5$, has $2$-positive
curvature operator, then volume-normalized Ricci flow converges to a unique
constant-positive-curvature metric.
\end{theorem}

\begin{theorem}[Ricci-flat metrics and Perelman's $\lambda$ (Theorem 35.41)]
\label{thm:chow-iv-ricci-flat-lambda-stability}
\source{chow-et-al-ricci-flow-part-iv:page-0326}
\uses{prop:chow-iv-perelman-lambda-second-variation}
\dcref{ch35:35.41}
\group{Alternative Stability Methods}


If a compact Ricci-flat metric $\hat g$ is a local maximizer of $\lambda$, then
every sufficiently close Ricci flow exists for all time and converges, modulo
diffeomorphisms, to a Ricci-flat metric near $\hat g$.
\end{theorem}

\begin{theorem}[Instability of non-maximizers (Theorem 35.42)]
\label{thm:chow-iv-ricci-flat-lambda-instability}
\source{chow-et-al-ricci-flow-part-iv:page-0327}
\dcref{ch35:35.42}
\group{Alternative Stability Methods}

If a compact Ricci-flat metric $\hat g$ is not a local maximizer of $\lambda$,
there exists a nontrivial ancient Ricci flow converging modulo diffeomorphisms
to $\hat g$ as $t\to-\infty$.
\end{theorem}

\subsection{Noncompact flat and negatively curved spaces}

For complete metrics on $\mathbb R^2$, Wu obtains immortal flows approaching a
soliton.  Schn\"urer--Schulze--Simon and Koch--Lamm prove global convergence of
small asymptotically Euclidean perturbations under coupled Ricci--harmonic or
Ricci--DeTurck flows.  Ye proves stability on compact negatively curved
manifolds when the operator
\[
 Lh_{ij}=\Delta_Lh_{ij}+\left(R^k{}_j+\tfrac1nR\,\delta^k_j\right)h_{ik}
\]
has strictly negative spectrum.  Li--Yin, Schn\"urer--Schulze--Simon, and Bamler
give corresponding stability results for noncompact hyperbolic and symmetric
spaces under decay, boundedness, or finite-volume hypotheses.

\begin{exercise}[Expanding soliton rescaling (Exercise 35.4)]
\label{ex:chow-iv-expanding-soliton-rescaling}
\dcref{ch35:35.4}
\group{Linear Stability}
\source{chow-et-al-ricci-flow-part-iv:page-0304}
Verify the expanding rescaling above by replacing $T-t$ with $t-T_0$ in the
shrinking-soliton calculation.
\end{exercise}

\begin{exercise}[Local form of the Ricci--DeTurck operator (Exercise 35.2)]
\label{ex:chow-iv-local-ricci-deturck}
\dcref{ch35:35.2}
\group{Linear Stability}
\source{chow-et-al-ricci-flow-part-iv:page-0303}
Show that
\begin{equation*}
 A_{\hat g}(g)_{ij}=a(x,\hat g,g)_{ij}^{k\ell pq}\partial_p\partial_qg_{k\ell}
 +b(x,\hat g,\partial\hat g,g)_{ij}^{k\ell p}\partial_pg_{k\ell}
 +c(x,\hat g,\partial\hat g,\partial^2\hat g)_{ij}^{k\ell}g_{k\ell},
 \tag{35.6}
\end{equation*}
with smooth dependence on $x$ and analytic dependence on the remaining
arguments.
\end{exercise}
